\section{Hardware}

\subsection{Protocol Hardening}
The redesign (Chapter 7) controls where each protocol may flow. 
This chapter defines how each protocol is secured on the device and conduit level.

\subsubsection{Profinet}
Profinet has no authentication or encryption features, so it cannot be hardened at the 
protocol level. Instead it is secured by isolation.
\begin{itemize}
    \item Layer 2 only - Profinet stays inside each OT Production station (VLAN 101 - 106).
        It is never routed, and never crosses through firewall. Even if routing would be 
        executed (through the attack) router would drop the communication.
    \item No inter-station Profinet communication - Each PLC talks Profinet (with I/O, HMI) only 
        within own VLAN.
    \item DCP Set disabled - desible DCP Set through TIA Portal for devices that support this feature.
        This prevents rename/reconfigure attacks through DCP.
\end{itemize}

\subsubsection{OPC UA}
OPC UA is the only routed OT data protocol. Hardening is specified in SR-03. \textit{(see Sercurity Requirements 6.1 TODO)}
\begin{itemize}
    \item Security mode Sign\&Encrypt - Security policy Basic256Sha256. Mode 'none' is disabled
        on all PLC and Festo Controllers servers.
    \item Certificate authentication - Clients (OPC UA proxy, Cobot/MLvisionPC) authenticate with application
        certificates. Annonymous and username/password sessions are disabled.
    \item Scoped access - The Cobot/MLvisionPC can reach only PLC2 on Station 2, and no other PLC's.
\end{itemize}

\subsubsection{MQTT}
MQTT is used by the Festo CXV device for telemetry. It's hardening is specified in SR-04 \textit{(see Sercurity Requirements 6.1 TODO)}
\begin{itemize}
    \item MQTT broker in IDMZ - The MQTT CXV device publishes only to a dedicated MQTT broker placed in IDMZ. IT nodes 
        can only subscribe to the broker. No direct OT-IT MQTT path exists.
    \item TLS 1.2+ - TLS encryption enforced on all MQTT connections.
    \item Client certificates - MQTT CXV device must authenticate with certificate for publish connection to the broker. No annonymous connections allowed.
    \item Topic ACL's - THe MQTT CXV device may publish only to it's own topics. IT nodes (MES) cannot publish to the OT topics.
\end{itemize}

\subsubsection{HTTP/Web Interfaces}
Some OT devices need to expose plaintext HTTP web interfaces. Hardening is specified in SR-05 \textit{(see Security Requirements 6.1 TODO)}

\begin{itemize}
    \item Disabled by default - All HTTP servers are disabled unless strictly needed for operational functionality.
    \item Reverse proxy - Where web access is needed (e.g. MES reading the PLC alarm table), the MES connects to the Nginx reverse proxy in the IDMZ over HTTPS/TLS 1.2+. Nginx terminates the connection, fetches the OT-side resource, and returns it. Plaintext HTTP stays inside the OT zone.
    \item Access restriction - Ngnix authenticate requests, logs them and blocks any OT server that should not be reached from IT zone.
        Firewall rules restrict web access to the Nginx TCP/IP port only.
\end{itemize}

\subsubsection{S7comm+ Engineering Access}
S7comm+ connection is used only for programming PLC's and HMI's. Hardening is sepcified in SR-06 and SR-07 TODO link to SR's
\begin{itemize}
    \item Single source - S7comm+/TCP102 to any PLC or HMI is allowed only from the dedicated Engineering Workstation in VLAN40.
    \item Conduit controll - Engineering Worksation S7comm+ connection initialization passes through the firewall, enforcing restriction of Engineering Station to PLC's and HMI's only.
    \item User Authentication - PLC's and HMI's access requires individual login/password authentication with role separation (implemented in TIA portal).
\end{itemize}



\subsection{Security Hardware}
This chapter lists the hardware required to implement the redesigned architecture (Chapter 7). 

\subsubsection{Industrial Firewall}

\renewcommand{\arraystretch}{1.3} % 30% more row height
\begin{tabularx}{\textwidth}{>{\bfseries}p{2.5cm} X}
    \toprule
    Model & Fortinet FortiGate Rugged 70G \\
    \midrule
    Function & Central Layer 3 firewall and inter-VLAN router. All traffic between VLANs passes through it - under default-deny polict. 
    Can perform stateful filtering and packet inspection of industrial protocols (S7comm+, OPC UA, MQTT, HTTP/S). Implements four logical firewalls 
    (Edge, IT–IDMZ, IDMZ–OT Cell, OT Cell–OT Production) as separate rule sets between VLAN interfaces. \\
    \midrule
    Placement & Network core. Has one tagged VLAN interface (.254 gateway) per VLAN. Has single 802.1Q trunk to the managed switch. No VLAN is bridged on the switch, so the firewall is the only inter-VLAN path. \\
    \midrule
\end{tabularx}

\subsubsection{Managed Switch}
\renewcommand{\arraystretch}{1.3} % 30% more row height
\begin{tabularx}{\textwidth}{>{\bfseries}p{2.5cm} X}
    \toprule
    Model & Siemens SCALANCE XC-216  \\
    \midrule
    Function & VLAN-aware Layer 2 switching with IEEE 802.1Q tagging. Provides per-station access ports (untagged, one VLAN each) and the trunk uplink to the firewall. \\
    \midrule
    Placement & Replace the existing unmanaged SCALANCE XB008 switches at each station. Interconnected by 802.1Q trunks; one trunk reaches the firewall. \\
    \midrule
\end{tabularx}

\subsubsection{Engineering Station}
\renewcommand{\arraystretch}{1.3} % 30% more row height
\begin{tabularx}{\textwidth}{>{\bfseries}p{2.5cm} X}
    \toprule
    Model & Siemens SIMATIC IPC627E\\
    \midrule
    Function & Dedicated managed station running TIA Portal for PLC and HMI programming/configuration. Only host permitted S7comm+ to PLCs/HMIs.\\
    \midrule
    Placement & OT Cell zone, isolated in its own VLAN (40). Reaches OT Production only through the firewall on TCP 102.\\
    \midrule
\end{tabularx}

\subsubsection{IDMZ proxy server}
\renewcommand{\arraystretch}{1.3} % 30% more row height
\begin{tabularx}{\textwidth}{>{\bfseries}p{2.5cm} X}
    \toprule
    Model & Industrial rack server or SIMATIC IPC. Hardened Linux.\\
    \midrule
    Function & Hosts the IDMZ services — OPC UA proxy, Nginx reverse proxy, MQTT broker, NTP master, log forwarder, backup staging. \\
    \midrule
    Placement & Placed inside IDMZ zone - on VLAN20. \\
    \midrule
\end{tabularx}

\subsubsection{Syslog \& Backup Storage}
\renewcommand{\arraystretch}{1.3} % 30% more row height
\begin{tabularx}{\textwidth}{>{\bfseries}p{2.5cm} X}
    \toprule
    Model & Standard server, Linux OS. \\
    \midrule
    Function & Central Syslog collector for firewall, switches, proxies and broker; offline configuration backups (after every change) and daily MES database backups over SFTP.\\
    \midrule
    Placement & IT zone, VLAN10\\
    \midrule
\end{tabularx}

